%% JORS (Journal of Open Research Software) submission rendering of Paper 1B.
%%
%% CANONICAL SOURCE: ../paper1b_namaster_proof.tex  (v2B.0.15, 2026-07-24)
%% This file is a SECOND RENDERING of that manuscript, mapped onto the JORS
%% Software Metapaper template (jors.cls + jors_template.tex, distributed as
%% JORS_Template.zip from the "LaTex Template" link on
%% https://openresearchsoftware.metajnl.com/about/submissions). It exists
%% because the JORS "Submission Preparation Checklist" (fetched 2026-07-24)
%% states, as one of the items "All submissions must meet": "The submission
%% conforms to the article template. For LaTex submissions, a PDF should be
%% provided along with the original LaTex file(s)."
%%
%% The canonical manuscript is NOT restructured: its v2B.0.13 bytes are
%% permanently archived under Zenodo DOI 10.5281/zenodo.21481842 and it feeds
%% the site, the PDF mirrors, and the arXiv tarball lineage. Two renderings of
%% one paper is expected — the journal imposes its own format.
%%
%% CONTENT RULE FOR THIS FILE: prose is carried VERBATIM from the canonical
%% manuscript. This is a reformat, never a content edit. Every honesty
%% disclosure of the canonical (the 41-test / 39-run + 2-monorepo-coupled-skip
%% contract, the macOS-untested label, the non-self-contained 1.41e-18 scalar
%% caveat, the "not measurements / not detection significances" framing, the
%% receipt-not-signature caveat, the non-affiliation-with-NaMaster statement,
%% and the AI Usage Disclosure) is preserved.
%%
%% SECTION MAPPING (canonical -> JORS skeleton):
%%   \title                     -> (1) Overview / Title
%%   \author                    -> Paper Authors
%%   Sec. Author Contributions  -> Paper Author Roles and Affiliations
%%   \begin{abstract}           -> Abstract
%%   Keywords paragraph         -> Keywords
%%   Sec. Introduction          -> Introduction
%%   Sec. Statement of Need     -> Introduction / Statement of need (subsection)
%%   Sec. Implementation & Arch -> Implementation and architecture
%%   Sec. Exact-Window Inference-> Implementation and architecture / subsection
%%   Sec. Content Validation    -> Implementation and architecture / subsection
%%   Sec. Quality Control       -> Quality control
%%   Sec. Worked Examples       -> Quality control / Worked examples (subsection)
%%   Sec. Availability (paras)  -> (2) Availability / Operating system,
%%                                 Programming language, Additional system
%%                                 requirements, Dependencies, Software
%%                                 location (Archive + Code repository),
%%                                 Validation artifacts
%%   (derived from repository)  -> (2) Availability / List of contributors
%%   (derived from repository)  -> (2) Availability / Language
%%   Sec. Reuse Potential       -> (3) Reuse potential
%%   Sec. Limitations           -> (3) Reuse potential / Limitations (subsection)
%%   Sec.* Acknowledgements     -> Acknowledgements
%%   Funding Statement para     -> Funding statement
%%   Competing Interests para   -> Competing interests
%%   Sec. AI Usage Disclosure   -> AI usage disclosure (additional section,
%%                                 placed with the other disclosure sections;
%%                                 no mandated section removed or reordered)
%%   thebibliography            -> References (carried verbatim, inline, so the
%%                                 bundle is self-contained and needs no bibtex
%%                                 pass and no citation renumbering)

\PassOptionsToPackage{hyphens}{url}
\documentclass{jors}

\usepackage{amsmath,amssymb}
\usepackage[T1]{fontenc}
\usepackage[utf8]{inputenc}
\usepackage{microtype}
\usepackage{seqsplit}
%% xcolor, enumitem, hyperref, fancyhdr, titlesec and sectsty are loaded by jors.cls

\hypersetup{
  colorlinks=true,
  linkcolor=blue!55!black,
  citecolor=blue!55!black,
  urlcolor=blue!55!black
}

%% jors.cls sets \sectionfont to 13pt but leaves \subsection at the 12pt-article
%% default (\Large, ~14.4pt), which would render subsections LARGER than the
%% mandated top-level sections. Restore a correct hierarchy without touching
%% jors.cls itself.
\subsectionfont{\fontsize{11.5}{13.5}\selectfont}

\newcommand{\software}{\texttt{namaster-proof}}
\newcommand{\Cl}{C_\ell}
\newcommand{\paperVersion}{v2B.0.15}
\newcommand{\paperTimestamp}{2026-07-24 12:00 PT}

%% Set the header information (per jors_template.tex)
\pagestyle{fancy}
\definecolor{mygray}{gray}{0.6}
\setlength{\headheight}{15pt}
\renewcommand\headrule{}
%% fancyhdr's default `fancy` style puts \leftmark on the left of the header,
%% which makes thebibliography's "REFERENCES" mark overprint the right-hand
%% note on the reference pages. Blank the left/centre header slots.
\lhead{}
\chead{}
\rhead{\textcolor{gray}{\footnotesize JORS rendering of \software{} metapaper
\paperVersion{} (\paperTimestamp)}}

\begin{document}

\rule{\textwidth}{1pt}

\section*{(1) Overview}

\vspace{0.5cm}

\section*{Title}

\software: Exact pseudo-\(\Cl\) window inference and content-bound validation
for reproducible spin-2 analyses

\section*{Paper Authors}

1. Golden, Houston

\section*{Paper Author Roles and Affiliations}

1. Houston Golden --- Independent researcher (unaffiliated), Los Angeles,
California, USA. ORCID:
\href{https://orcid.org/0009-0008-5616-5994}{0009-0008-5616-5994}.
Correspondence: \href{mailto:houston@hubify.com}{houston@hubify.com}.

Author roles: Houston Golden: conceptualization, methodology, software,
validation, formal analysis, investigation, data curation,
writing---original draft, writing---review and editing, visualization, project
administration, and funding acquisition. Correspondence metadata remain
author-supplied submission metadata and are not inferred by the software
release process.

\section*{Abstract}

\software{} is a focused Python verification layer for two error-prone steps
in cut-sky spin-2 analyses.  First, it evaluates a uniformly rotated
\(EE,EB,BE,BB\) spectrum through the complete NaMaster bandpower-window
operator, avoiding replacement of the operator by bin-centre or
effective-multipole templates.  Second, it writes JSON results and
content-bound sidecar receipts with atomic per-file replacement and fails
closed when result bytes or caller-asserted execution metadata change.  The package also supplies explicit
multipole-support contracts, fixed-grid rotation-angle recovery, command-line
receipt verification, and compatibility tests against the production helpers
from which it was extracted.  The software is intended for method validation
and reproducibility checks; it is not a sky-analysis pipeline, foreground
model, or cosmological inference engine.

\section*{Keywords}

Python; cosmology; pseudo-\(\Cl\); NaMaster; reproducibility; provenance.

\section*{Introduction}

Pseudo-\(\Cl\) estimators make masked spin-2 analyses computationally
practical by relating full-sky spectra to coupled and decoupled
bandpowers~\cite{Hivon2002,Alonso2019}.  Reproducible use of these estimators
requires more than recording a best-fit number: the same window operator must
be used in simulation and inference, harmonic limits must agree across fields
and bins, and intermediate results must remain bound to their execution
metadata.

\software{} packages these checks behind a small public API.  Its numerical
module constructs exact windowed responses for uniform polarization rotation,
evaluates them at one or many angles, recovers an angle on a declared grid,
and compares the window contraction with NaMaster's couple--decouple operator.
Its provenance module publishes a JSON result and a SHA-256-bound receipt using
atomic replacement for each file, then verifies the pair and caller-asserted
metadata.  A third module defines
deterministic multipole and binning contracts.  NumPy is the only required
dependency; a NaMaster workspace is supplied by the user when exercising the
physical window operator.

\subsection*{Statement of need}

For two spin-2 fields, NaMaster exposes a bandpower-window tensor with spectrum
ordering \([EE,EB,BE,BB]\).  A common shortcut is to evaluate a theory
spectrum at a representative multipole for each bin.  That shortcut is not
generally identical to applying the complete mode-coupling and binning
operator, and a mismatch between the simulation and fitting operators can
appear as a parameter-recovery bias.

A separate reproducibility risk arises in long or sharded calculations.
Ordinary filenames do not prove that a summary corresponds to the intended
seeds, realization count, configuration, or result bytes.  A process can also
terminate between writing a result and its metadata, leaving a plausible but
incomplete artifact.  These problems are usually addressed with analysis-local
helpers.  \software{} provides a reusable, tested implementation with explicit
failure behavior, while remaining narrow enough to audit.

\section*{Implementation and architecture}

The package is divided into three modules:

\begin{description}[leftmargin=2.8cm,style=nextline]
\item[\texttt{windows}]
  Rotation of \(EE,EB,BE,BB\) spectra, exact contraction with a
  \nolinkurl{get_bandpower_windows()} tensor, fixed-grid angle recovery, and
  an equivalence check against \nolinkurl{couple_cell()} followed by
  \nolinkurl{decouple_cell()}.
\item[Multipoles]
  Validation of a common harmonic cutoff and construction of integer
  bandpower edges whose final exclusive edge is \(\ell_{\max}+1\).
\item[\texttt{receipts}]
  Atomic JSON publication, streaming SHA-256 calculation, sidecar-receipt
  construction, and fail-closed content and metadata validation.
\end{description}

The functions accept array-like inputs and return NumPy arrays or plain Python
dictionaries.  NaMaster is deliberately not an installation dependency:
window functions require a workspace exposing
\nolinkurl{get_bandpower_windows()}, \nolinkurl{couple_cell(cls)}, and
\nolinkurl{decouple_cell(coupled)}. The first must return a finite
\([4,n_b,4,n_\ell]\) tensor and the couple--decouple route must return a finite
\([4,n_b]\) array. This small protocol makes the algebra testable with synthetic
linear operators and permits users to pin PyMaster independently.

\subsection*{Exact-window inference}

For an initially vanishing \(EB\) spectrum and a uniform rotation angle
\(\beta\), the rotated spectra can be decomposed into angle-independent,
\(\cos(4\beta)\), and \(\sin(4\beta)\) components.  In particular,
\begin{align}
 C_\ell^{EE}(\beta)
   &= \tfrac12(C_\ell^{EE}+C_\ell^{BB})
    + \tfrac12(C_\ell^{EE}-C_\ell^{BB})\cos(4\beta),\\
 C_\ell^{BB}(\beta)
   &= \tfrac12(C_\ell^{EE}+C_\ell^{BB})
    - \tfrac12(C_\ell^{EE}-C_\ell^{BB})\cos(4\beta),\\
 C_\ell^{EB}(\beta)=C_\ell^{BE}(\beta)
   &= \tfrac12(C_\ell^{EE}-C_\ell^{BB})\sin(4\beta).
\end{align}
Let \(W_{b\ell}^{ij}\) be the workspace tensor mapping input spectrum \(j\)
to output spectrum \(i\) in band \(b\).  The package precontracts each of the
three components,
\begin{equation}
 R_{b}^{i,k}=\sum_{j,\ell}W_{b\ell}^{ij}C_\ell^{j,k},
\end{equation}
and subsequently evaluates
\begin{equation}
 C_b^i(\beta)=R_b^{i,0}
 +\cos(4\beta)R_b^{i,c}+\sin(4\beta)R_b^{i,s}.
\end{equation}
This retains the full window tensor while making repeated grid evaluation
inexpensive.  The recovery function minimizes a declared weighted or
unweighted squared residual over a user-provided grid; its default is 2001
points spanning \(-1^\circ\) to \(+1^\circ\).

\subsection*{Content validation}

\texttt{publish\_json} serializes a mapping with sorted keys, derives the
receipt digest and byte count from that immutable serialized snapshot, writes and
flushes a process-specific temporary file, atomically replaces the result,
and synchronizes the containing directory.  It then separately publishes a sidecar receipt
containing schema version, result filename, byte count, SHA-256 digest, and
caller-supplied execution metadata.  The content-binding fields are protected
from metadata override.

\texttt{verify\_json\_receipt} requires both files, reads each through one file
handle, parses those immutable byte snapshots as JSON objects, and recomputes
every protected field from the same result bytes it returns.
\texttt{validate\_json\_receipt}
additionally checks caller-asserted metadata such as a suite identifier, seed range,
or realization count using recursive JSON-type-strict equality, so booleans,
integers, and floating-point numbers are not interchangeable. Missing files,
malformed objects, changed bytes,
changes to protected receipt fields, or asserted-contract mismatches raise an exception. The two
files are not one filesystem transaction, and coordinated replacement of both
requires an external receipt anchor or trusted expected metadata to detect. The
\texttt{namaster-proof verify} and \texttt{namaster-proof validate} commands
expose the same behavior to shell workflows and return a nonzero status on
failure.

\section*{Quality control}

Version 0.1.7 contains 41 automated tests covering numerical behavior, invalid
inputs, the CLI, receipt mutation, protected fields, and compatibility with
the original production helpers.  All 41 run in a monorepo checkout; in a
standalone package install 39 run and the 2 replay-equivalence tests skip
cleanly, because those two tests replay against production helper modules that
reside elsewhere in the monorepo rather than inside the package.  The window suite uses an exactly represented
synthetic linear workspace to compare the precontracted response with the
couple--decouple route and to recover an injected grid angle.  Rotation tests
check conservation of total \(EE+BB\) auto-power.  Receipt tests mutate result
bytes and receipt digests independently and require rejection.

Compatibility tests load the pre-package production modules and compare
windowed bandpowers at negative, zero, and positive angles with zero requested
absolute or relative tolerance.  They also compare algebraic spectrum rotation,
bandpower edges, and harmonic keyword contracts.  In the associated physical
production artifact, a workspace of shape
\([4,20,4,1025]\) gave a maximum absolute difference of
\(1.41\times10^{-18}\) between direct window contraction and the
couple--decouple operator. This recorded execution check used IEEE-754
double-precision arrays and the production driver's \(10^{-10}\) acceptance
threshold; the package's separate synthetic legacy-helper comparison uses
zero requested absolute and relative tolerance. Because
the original workspace tensor was not retained, the scalar is not a
self-contained reproducibility claim or a universal error bound.
The tensor is nonetheless deterministically regenerable from committed,
RNG-free code---the mask is a pure function of \(N_{\rm side}\) and fixed
latitude cuts and the \(10^{-10}\)-gated equivalence check is committed---so
the \nolinkurl{examples/rebuild_workspace_check.py} script rebuilds the
\([4,20,4,1025]\) workspace and re-verifies \(\max|\Delta|<10^{-10}\) whenever
PyMaster is installed (only the last digits of the \(\sim10^{-18}\) value are
platform-dependent). The full suite runs with
\texttt{python -m pytest packages/namaster-proof/tests} from the repository
root (or \texttt{pytest} from within the installed package's test extra).

\subsection*{Worked examples}

\paragraph{Minimal synthetic operator.}
The documentation constructs an identity-like workspace, builds a response
from synthetic \(EE\) and \(BB\) arrays, evaluates a known rotation, and checks
operator equivalence.  This example has no dependency on a paper-specific
dataset and is suitable for installation tests.
The minimal call sequence is
\texttt{response = build\_rotation\_response(workspace, ee, bb)} followed by
\texttt{windowed\_bandpowers(response, beta\_rad=np.deg2rad(0.25))}; the complete executable
example also calls \texttt{validate\_window\_equivalence}.

\paragraph{Real PyMaster integration.}
The optional \nolinkurl{examples/pymaster_integration.py} example constructs a
real PyMaster workspace, verifies the exact operator equivalence, recovers a
declared grid injection, and records the result beside an effective-multipole
shortcut comparison. It requires separately installed PyMaster and healpy.

\paragraph{Synthetic CMB recovery campaign.}
The production use case employed CAMB 1.6.6 raw \(EE/BB\) spectra, PyMaster
2.6, \(N_{\rm side}=512\), \(\ell_{\max}=1024\), and deterministic seeds.
For each of two nonzero injected angles, a 500-realization, foreground-free
synthetic run recovered the injection from the mean bandpowers on the declared
\(0.001^\circ\) grid: \(0.270^\circ\mapsto0.270^\circ\) and
\(0.342^\circ\mapsto0.342^\circ\).  The null injection, also run with 500
realizations, recovered \(0.000^\circ\).  These are software-recovery checks under the stated
simulation contract, not measurements, detection significances, or evidence
for a physical birefringence model.

\paragraph{Sharded result validation.}
Production shards record fields such as suite, realization count, and seed
range alongside the content digest.  Aggregation can therefore reject a shard
whose bytes were altered or whose declared range does not match the requested
campaign, rather than silently combining it with valid results.

\section*{(2) Availability}
\vspace{0.5cm}

\section*{Operating system}

The package supports Linux and Windows where Python and NumPy are available;
CI covers Linux with Python 3.10--3.13 and Windows with Python 3.12. Directory
metadata are synchronized on POSIX systems after atomic replacement. macOS is
expected to work as a POSIX platform wherever Python and NumPy are available,
but is not currently exercised in continuous integration and is therefore
listed as untested.

\section*{Programming language}

Python. Version 0.1.7 requires Python 3.10 or later.

\section*{Additional system requirements}

The package has no special hardware requirements: it is pure Python with a
single NumPy runtime dependency, requires no GPU, and its full test suite
completes in under one second with a negligible memory footprint on a commodity
CPU. The optional PyMaster integration example (\(N_{\rm side}=16\)) completes
in a few seconds, whereas the retained physical validation campaign
(\(N_{\rm side}=512\), \(\ell_{\max}=1024\); 500 realizations at each of three
injected angles) recorded a wall time of order \(7\times10^{2}\,\mathrm{s}\)
with eight parallel realization workers on a commodity multi-core CPU
(correspondingly of order a couple of hours single-worker). These campaign
timings are environment-dependent; both workloads depend on the user-supplied
PyMaster and healpy stacks and correspondingly more memory, but lie outside the
self-contained packaged code path.

\section*{Dependencies}

NumPy 1.24 or later is the only required runtime dependency. PyMaster
is optional and is supplied by the user for physical NaMaster workspaces;
the retained physical validation used PyMaster 2.6 and healpy 1.19.0.

\section*{List of contributors}

Houston Golden (Independent researcher, Los Angeles, California, USA) is the
sole author and contributor of the software --- design, implementation, tests,
examples, documentation, and packaging --- consistent with the package's
\texttt{CITATION.cff} and \texttt{codemeta.json}, which each list a single
author. AI coding and review agents assisted under the author's supervision;
see the \emph{AI usage disclosure} section below for the scope of that
assistance and how its outputs were verified.

\section*{Software location:}

{\bf Archive}

\begin{description}[noitemsep,topsep=0pt]
	\item[Name:] \software{} 0.1.7 --- commit-pinned archive of the
	  \nolinkurl{packages/namaster-proof} source tree (repository commit
	  \texttt{\seqsplit{0a587b583f8e86c4ce1ee4a20526fcdcd8035fe6}}), together
	  with its license, citation metadata, and checksums.
	\item[Persistent identifier:]
	  \href{https://doi.org/10.5281/zenodo.21481753}{doi:10.5281/zenodo.21481753}
	\item[Licence:] MIT License.
	\item[Publisher:] Zenodo.
	\item[Version published:] 0.1.7
	\item[Date published:] 21 July 2026
\end{description}

The manuscript itself, with its exact source, arXiv bundle, and provenance
manifest, is separately archived under
\href{https://doi.org/10.5281/zenodo.21481842}{doi:10.5281/zenodo.21481842}
(CC-BY-4.0, deposited July 21, 2026).

{\bf Code repository}

\begingroup\sloppy
\begin{description}[noitemsep,topsep=0pt]
	\item[Name:] \nolinkurl{Hubify-Projects/bigbounce} ---
	  \nolinkurl{packages/namaster-proof}
	\item[Identifier:]
	  \url{https://github.com/Hubify-Projects/bigbounce/tree/main/packages/namaster-proof}
	\item[Licence:] MIT License.
	\item[Date published:] 16 July 2026
\end{description}
\par\endgroup

\begingroup\sloppy
The software is released under the MIT License.
The source, issue tracker, installation instructions, API documentation, tests,
and examples are available at
\url{https://github.com/Hubify-Projects/bigbounce/tree/main/packages/namaster-proof}.
From a checkout of that repository the package installs locally with
\texttt{python -m pip install ./packages/namaster-proof} (append
\texttt{-e '.[test]'} from within the package directory for the test extra).
\texttt{namaster-proof} is an independent downstream verification layer and is
not an official release of, or affiliated with, the NaMaster project.
Citation metadata are supplied in \texttt{CITATION.cff} and machine-readable
metadata in \texttt{codemeta.json}. The package directory was first published
to the repository on 2026-07-16.
\par\endgroup

\subsection*{Validation artifacts}

The corrected physical example is bound to the
\href{https://github.com/Hubify-Projects/bigbounce/blob/main/reproducibility/p1_namaster_500mc/results/physical_spectrum_v2/summary.json}
{repository summary artifact} (SHA-256
\texttt{\seqsplit{745b0a2f060773ce69c005ea84b74b305ec26a85f6aaafe58f0b3244b7f39914}})
and the
\href{https://github.com/Hubify-Projects/bigbounce/blob/main/reproducibility/p1_namaster_500mc/results/physical_spectrum_v2/bandpowers.npz}
{repository bandpower artifact} (SHA-256
\texttt{\seqsplit{b00f850e338007caea6af76f4e9305ab6b54a68e6799efd450bc76f1c325f331}}).

\section*{Language}

English. The repository, software documentation, source comments, and all
supporting files are written in English; the software itself is implemented in
Python.

\section*{(3) Reuse potential}

Researchers can reuse the window layer to test uniform spin-2 rotation models
against any workspace exposing the small NaMaster protocol, reuse the multipole
contracts to prevent field/bin support drift, or use the receipt layer for
sharded JSON calculations outside cosmology. The receipt primitive remains in
the same package because it authenticates the exact validation artifacts
produced by the window workflows; it is a small reusable module rather than a
claim that the package is a general provenance framework. Extensions can add other
angle-dependent spectrum bases, signed receipt anchors, or schema migrations.
Questions, bug reports, and contributions are handled through the public
repository issue tracker and pull-request workflow.

\subsection*{Limitations}

\software{} does not create masks or maps, estimate covariance matrices,
separate cosmic rotation from instrumental polarization calibration, model
foregrounds, choose a likelihood, or infer cosmological parameters.  Uniform
rotation and initially vanishing \(EB\) are assumptions of the optimized
three-component response; the general algebraic rotation helper accepts all
four input spectra, but spatially varying rotation requires another model.
Fixed-grid recovery inherits the range and resolution selected by the caller.
Receipt validation establishes file integrity and agreement with caller-asserted
metadata, not the scientific correctness of that metadata. SHA-256 sidecars
are content bindings, not digital signatures, an authorization system, or
protection against coordinated replacement without an external anchor.
Schema version 1 has no migration requirement yet; a future incompatible
schema will require an explicit reader or conversion path.

\section*{Acknowledgements}

The exact-window operations build on NaMaster~\cite{Alonso2019} and the MASTER
pseudo-\(\Cl\) method~\cite{Hivon2002}.  The production example used
CAMB~\cite{Lewis2000} and HEALPix~\cite{Gorski2005}.

\section*{Funding statement}

No external funding supported this software release.

\section*{Competing interests}

The author declares no competing interests.

\section*{AI usage disclosure}

AI coding and review agents assisted with implementation review, test
generation, documentation, and manuscript editing. Their outputs were checked
against committed source and retained numerical artifacts; executable claims
were accepted only after local tests or direct recomputation.

\footnotesize
\begin{thebibliography}{9}

\bibitem{Hivon2002}
E.~Hivon, K.~M. G\'orski, C.~B. Netterfield, B.~P. Crill, S.~Prunet, and
F.~Hansen, ``MASTER of the CMB anisotropy power spectrum:
A fast method for statistical analysis of large and complex cosmic microwave
background data sets,'' \emph{Astrophysical Journal} 567, 2 (2002),
\href{https://doi.org/10.1086/338126}{doi:10.1086/338126}.

\bibitem{Alonso2019}
D.~Alonso, J.~Sanchez, and A.~Slosar, ``A unified pseudo-\(C_\ell\) framework,''
\emph{Monthly Notices of the Royal Astronomical Society} 484, 4127--4151
(2019), \href{https://doi.org/10.1093/mnras/stz093}{doi:10.1093/mnras/stz093}.

\bibitem{Lewis2000}
A.~Lewis, A.~Challinor, and A.~Lasenby, ``Efficient computation of cosmic
microwave background anisotropies in closed Friedmann--Robertson--Walker
models,'' \emph{Astrophysical Journal} 538, 473--476 (2000),
\href{https://doi.org/10.1086/309179}{doi:10.1086/309179}.

\bibitem{Gorski2005}
K.~M. G\'orski, E.~Hivon, A.~J. Banday, B.~D. Wandelt, F.~K. Hansen,
M.~Reinecke, and M.~Bartelmann, ``HEALPix: A framework for high-resolution
discretization and fast analysis of data distributed on the sphere,''
\emph{Astrophysical Journal} 622, 759--771 (2005),
\href{https://doi.org/10.1086/427976}{doi:10.1086/427976}.

\end{thebibliography}
\normalsize

\vspace{2cm}

\rule{\textwidth}{1pt}

{ \bf Copyright Notice} \\
Authors who publish with this journal agree to the following terms: \\

Authors retain copyright and grant the journal right of first publication with the work simultaneously licensed under a \href{http://creativecommons.org/licenses/by/3.0/}{Creative Commons Attribution License} that allows others to share the work with an acknowledgement of the work's authorship and initial publication in this journal. \\

Authors are able to enter into separate, additional contractual arrangements for the non-exclusive distribution of the journal's published version of the work (e.g., post it to an institutional repository or publish it in a book), with an acknowledgement of its initial publication in this journal. \\

By submitting this paper you agree to the terms of this Copyright Notice, which will apply to this submission if and when it is published by this journal.

\end{document}
